% TOP_PROBLEM: {"id": 30006951, "title": "Greenfield-Wallach-Katok Conjecture for Cohomology-Free Vector Fields", "category_id": 11, "category": {"id": 11, "name": "dynamical_systems", "display_name": "Dynamical Systems", "description": "Problems about long-term behavior of deterministic systems, Hamiltonian dynamics, and periodic orbits.", "slug": "dynamical-systems", "order_index": 11, "created_at": "2026-07-31T15:26:25.671Z"}, "status": "open"}
% FIRST_POSTED: 2026-09-27
% !TeX program = lualatex
% ATTEMPT_STATUS: unresolved
\documentclass[11pt]{article}
\usepackage[margin=25mm]{geometry}
\usepackage{fontspec}
\setmainfont{Latin Modern Roman}
\usepackage{amsmath,amssymb,mathtools}
\usepackage{unicode-math}
\setmathfont{Latin Modern Math}
\usepackage{xurl,hyperref}
\hypersetup{colorlinks=true,urlcolor=blue}
\setcounter{secnumdepth}{0}
\setlength{\parindent}{0pt}
\setlength{\parskip}{5pt}
\setlength{\emergencystretch}{3em}
\begin{document}
\section*{\texorpdfstring{Greenfield-Wallach-Katok Conjecture for Cohomology-Free Vector Fields}{Greenfield-Wallach-Katok Conjecture for Cohomology-Free Vector Fields}}\label{30006951}
\subsection{Definitions and mathematical statement}
Let $M$ be a closed connected orientable smooth $d$-manifold, $d\ge1$, and $X$ a smooth vector field. Write $Xu={\,\mathrm{d}} u(X)$. It is cohomology-free when
\[
 \forall v\in C^\infty(M,{\mathbb{R}})\quad
 \exists u\in C^\infty(M,{\mathbb{R}}),\ c\in{\mathbb{R}}:\quad Xu=v-c.
\]
Must there be a smooth diffeomorphism $h:M\to\mathbb T^d={\mathbb{R}}^d/{\mathbb{Z}}^d$ for which $h_*X$ is a constant vector $\alpha\in{\mathbb{R}}^d$ satisfying a Diophantine estimate
\[
 |k\cdot\alpha|\ge C\|k\|^{-\tau}
 \quad(0\ne k\in{\mathbb{Z}}^d),\qquad C>0,\quad\tau\ge0?
\]
The constants may depend on the vector field. The lower bound excludes resonances as well as excessively close rational approximations. The catalogue identifies this cohomological formulation with the globally hypoelliptic form; the displayed target uses smooth real functions throughout.
\par\smallskip\textit{Formulation and concept references:} \href{https://arxiv.org/pdf/0706.3981}{[S1]}.

\subsection{Short English statement}
Must a flow for which every smooth cohomological equation is solvable up to a constant be a Diophantine linear flow on a torus?
\subsection{Research attempt: invariant volume, small divisors, and the torus factor}
\textbf{Status.} The general classification is not established. The argument proves the classification for constant torus fields, constructs a smooth torus factor of dimension $b_1(M)$ for every field in the question, and obtains the full conclusion when $b_1(M)=d$. It leaves an explicit topological gap when $b_1(M)<d$.

\textbf{Source audit.} The retrieved primary manuscript [S1] is Giovanni Forni's arXiv:0706.3981v1, whose PDF title specifies dimension three. Its Theorem 4.3 gives a torus fibration reduction, and Section 5 treats the three-dimensional classification. The general-dimensional conclusion is not a consequence of the title alone. The reductions below give a self-contained route through several of the underlying ideas; they are established partial theory, not a claimed new solution. We work only with the smooth cohomological property stated in the question and need no unproved change of formulation to global hypoellipticity.

Let $\phi_t$ be the complete flow of $X$; compactness of $M$ ensures completeness. Write $b=b_1(M)=\dim H^1(M;\mathbb R)$.

\subsubsection{1. The invariant measure is smooth and unique}
Choose an arbitrary smooth positive volume form $\rho$ on the orientable manifold. Define its divergence by
\[
 \mathcal L_X\rho=(\operatorname{div}_\rho X)\rho.
\]
Apply the cohomology-free property to $v=\operatorname{div}_\rho X$ and obtain
\[
 Xu=\operatorname{div}_\rho X-c.
\]
The positive volume form $\eta=e^{-u}\rho$ then satisfies
\[
 \mathcal L_X\eta=(\operatorname{div}_\rho X-Xu)\eta=c\eta.
\]
Cartan's formula and Stokes's theorem on a closed manifold give
\[
 0=\int_M d(\iota_X\eta)=\int_M\mathcal L_X\eta
   =c\int_M\eta.
\]
Thus $c=0$. Normalizing $\eta$ produces a smooth invariant probability measure $\mu$ with positive density everywhere.

For any equation $Xu=v-c$, invariance gives $\int Xu\,d\mu=0$, and hence
\[
 c=\int_M v\,d\mu.                                      \tag{1}
\]
In particular the constant is uniquely determined, independently of the solution $u$. If $\nu$ is any other invariant probability measure, the same equation gives $\int v\,d\nu=c$ for every smooth $v$. Smooth functions are uniformly dense in continuous functions on a compact manifold; therefore $\nu=\mu$. This proves unique ergodicity, not merely existence of a preferred smooth invariant measure.

Every nonempty closed flow-invariant subset $K$ supports an invariant probability: average a point mass along an orbit in $K$, take a weakly convergent subsequence, and note that translating an averaging interval changes its integral against a bounded function by $O(1/T)$. Uniqueness forces this measure to be $\mu$, whose support is all of $M$. Hence $K=M$, and the flow is minimal.

There are no zeros of $X$, since a zero would support an invariant point mass distinct from smooth volume. In dimension $d\geq2$ there are also no periodic orbits: their normalized orbit measures are singular to smooth $d$-dimensional volume. The latter conclusion must not be asserted in dimension one, where a nonvanishing constant circle flow is periodic as a continuous-time flow and is nevertheless cohomology-free.

The argument also controls invariant distributions. If a continuous linear functional $D$ on $C^\infty(M)$ satisfies $D(Xf)=0$ for every smooth $f$, then the cohomological equation and (1) give
\[
 D(v)=D(1)\int_M v\,d\mu.                               \tag{2}
\]
Thus all invariant distributions are multiples of the smooth invariant volume functional. This is substantially stronger than a statement about invariant probabilities alone, but it still does not identify the topology of $M$.

\subsubsection{2. Constant torus fields: necessity and sufficiency of the arithmetic bound}
For $X_\alpha=\alpha\cdot\nabla$ on $\mathbb T^m$, use Fourier functions $e_k(x)=e^{2\pi i k\cdot x}$. Then
\[
 X_\alpha e_k=2\pi i(k\cdot\alpha)e_k.                   \tag{3}
\]
If $|k\cdot\alpha|\geq C\|k\|^{-\tau}$ for $k\ne0$, every smooth $v$ has the solution
\[
 c=\widehat v(0),\qquad \widehat u(0)=0,\qquad
 \widehat u(k)=\frac{\widehat v(k)}{2\pi i k\cdot\alpha}
 \quad(k\ne0).                                          \tag{4}
\]
The Fourier coefficients of a smooth function decay faster than any inverse power. The denominator in (4) loses only a fixed power, so the same is true for $u$. All derivatives of its Fourier series converge uniformly after choosing a sufficiently strong decay exponent. Real $v$ gives real $u$ by conjugate symmetry. More quantitatively,
\[
 \|u\|_{H^s}\leq C_s\|v-\widehat v(0)\|_{H^{s+\tau}}
                                                               \tag{5}
\]
up to the fixed equivalence between $\|k\|$ and $(1+\|k\|^2)^{1/2}$ on nonzero frequencies.

Conversely, a resonance $k\cdot\alpha=0$ with $k\ne0$ makes $v=\cos(2\pi k\cdot x)$ unsolvable: its nonzero $k$th Fourier coefficient cannot equal the $k$th coefficient of $X_\alpha u$, whatever constant is subtracted.

Suppose now that there are no resonances but no Diophantine lower bound. We may choose distinct frequencies, no two equal up to sign, with
\[
 \|k_j\|\geq2^j,\qquad
 0<|k_j\cdot\alpha|<\|k_j\|^{-j^2}.                     \tag{6}
\]
To justify choosing them successively outside any finite set, note that on a fixed finite set the nonzero values $|k\cdot\alpha|$ have a positive minimum after multiplication by any prescribed power. If arbitrarily strong small-divisor failures occurred only in that set, an appropriate positive constant would give a global bound for that power, contrary to the assumption. Taking the constant small enough excludes all previously chosen frequencies and all bounded frequencies.

Set $a_j=|k_j\cdot\alpha|^{1/2}$ and
\[
 v(x)=\sum_{j\geq1} a_j\cos(2\pi k_j\cdot x).
\]
For every fixed derivative order $q$, the terms are bounded by a constant times $\|k_j\|^{q-j^2/2}$, a summable sequence after discarding finitely many initial indices. Thus $v$ is smooth and has zero mean. But any solution of $X_\alpha u=v$ must satisfy
\[
 |\widehat u(k_j)|=\frac1{4\pi|k_j\cdot\alpha|^{1/2}}
 >\frac1{4\pi}\|k_j\|^{j^2/2}.                           \tag{7}
\]
These coefficients are unbounded, and indeed grow faster than every fixed polynomial. They cannot be the Fourier coefficients of a smooth function, or even of a distribution. This proves that constant torus fields are cohomology-free exactly under the displayed Diophantine condition.

In particular, merely requiring $k\cdot\alpha\ne0$ is insufficient. Rational independence gives minimal translation dynamics; it does not control division by the small quantities in (3).

\subsubsection{3. Constructing a torus factor from closed one-forms}
Assume first $b>0$. Choose smooth closed one-forms $\beta_1,\ldots,\beta_b$ representing an integral basis of the free part of $H^1(M;\mathbb Z)$. Their periods are integers. Solve
\[
 Xu_i=\beta_i(X)-c_i,\qquad
 \eta_i=\beta_i-du_i.
\]
Then $d\eta_i=0$, all periods remain integral, and $\eta_i(X)=c_i$. Choosing a basepoint and integrating along paths defines a well-defined smooth map
\[
 h:M\longrightarrow\mathbb T^b,\qquad
 h(x)=\left(\int_{x_0}^x\eta_1,\ldots,
                 \int_{x_0}^x\eta_b\right)\pmod{\mathbb Z^b}.
                                                               \tag{8}
\]
Different paths change each integral by an integer. Local primitives show smoothness and $dh=(\eta_1,\ldots,\eta_b)$. Hence, for $c=(c_1,\ldots,c_b)$,
\[
 h(\phi_t x)=h(x)+tc.                                   \tag{9}
\]

The vector $c$ has no integer resonance. If $0\ne k\in\mathbb Z^b$ had $k\cdot c=0$, the circle-valued function $k\cdot h$ would be flow invariant. Minimality forces every continuous invariant function to be constant. Differentiating this constant circle map locally gives
\[
 \sum_i k_i\eta_i=0.
\]
Its cohomology class is $\sum_i k_i[\beta_i]$, contradicting linear independence. The translation flow with velocity $c$ is therefore dense on $\mathbb T^b$. One way to see the last implication is that the closure of a one-parameter subgroup of a torus is a closed subgroup, and a proper such subgroup has a nontrivial annihilating integer character. That character would give a resonance.

The compact set $h(M)$ is invariant under that dense translation flow by (9), so $h(M)=\mathbb T^b$. We also need to rule out critical points; surjectivity alone does not do this. Differentiating (9) gives
\[
 dh_{\phi_t x}\circ d\phi_t(x)=dh_x,
\]
so the rank of $dh$ is constant along each orbit. For each integer $q$, the set $\{x:\operatorname{rank}dh_x\leq q\}$ is closed and invariant. Minimality makes it empty or all of $M$. Thus the rank is globally constant. If it were less than $b$, Sard's theorem would make $h(M)$ a set of measure zero in $\mathbb T^b$, contradicting surjectivity. Therefore $h$ is a submersion and
\[
 b_1(M)\leq d.                                         \tag{10}
\]
Compactness makes this a proper submersion, so it has compact fibers and locally trivial smooth fiber charts. If $b=0$, this construction is simply the map to a point and asserts no positive-dimensional factor.

\subsubsection{4. The factor is itself cohomology-free}
For clarity, the arithmetic conclusion about $c$ needs an argument beyond its nonresonance. Push the smooth invariant volume $\mu$ forward by $h$. This is a probability invariant under a dense translation flow, hence Haar probability $dy$ on $\mathbb T^b$. For a smooth $u$ on $M$, fiber integration defines the smooth function $Pu$ characterized by
\[
 h_*(u\mu)=(Pu)(y)\,dy.                                 \tag{11}
\]
Smoothness follows in submersion charts by integration over fiber variables, using a partition of unity; compactness permits finitely many such local integrations over each coordinate patch. In particular $P1=1$.

Invariance of $\mu$ and (9) imply
\[
 P(u\circ\phi_t)(y)=Pu(y+tc).
\]
This can be verified by integrating both sides against a smooth test function on the torus and changing variables $x\mapsto\phi_t x$. Differentiation at $t=0$ gives
\[
 P(Xu)=c\cdot\nabla(Pu).                                \tag{12}
\]
Given smooth $g$ on $\mathbb T^b$, solve $Xu=g\circ h-a$ on $M$. Applying (12) yields
\[
 c\cdot\nabla(Pu)=g-a.
\]
Thus the constant factor field is cohomology-free. Section 2 now proves a Diophantine estimate for $c$. This order of reasoning is necessary: nonresonance alone would not give that estimate.

If $b=d$, the submersion $h$ has discrete compact fibers, so it is a finite covering. A connected finite cover of $\mathbb R^d/\mathbb Z^d$ is $\mathbb R^d/\Lambda$ for a finite-index lattice $\Lambda\subset\mathbb Z^d$. In these covering coordinates the lifted field is the constant vector $c$, because $dh(X)=c$ and the covering map is locally the identity. A linear identification of $\Lambda$ with $\mathbb Z^d$ is a smooth diffeomorphism to the standard torus. The resulting constant field is cohomology-free by conjugacy invariance, and Section 2 gives its Diophantine condition. This proves the full target in the case $b_1(M)=d$.

\subsubsection{5. Remaining gap and checks on tempting shortcuts}
The first general obstruction is now specific: the preceding arguments do not prove $b_1(M)=d$. If $b<d$, the fibers in (8) have positive dimension, and if $b=0$ they give no torus information at all. Minimality does not make those fibers points: the flow carries one fiber to another and does not preserve an individual fiber when the base velocity is nonzero. Therefore restricting $X$ to a fiber and applying the conjecture inductively is not a legitimate operation.

Likewise, unique ergodicity and absence of periodic orbits alone are far weaker than solving every smooth cohomological equation. The small-divisor construction (6)--(7) supplies explicit uniquely ergodic torus flows for which smooth solvability fails. A proof using only unique ergodicity has lost an essential analytic part of the hypothesis.

The torus formulas can be checked on finite Fourier sums, but finite frequency tests cannot certify a Diophantine lower bound over all integer frequencies. Accordingly the accompanying diagnostic checks only Fourier inversion identities; it is not presented as a test for the conjecture. The invariant-volume construction, rank argument, and fiber-averaging argument are analytic proofs independent of computation.

The completed reduction gives a smooth Diophantine torus factor and proves the classification when the first Betti number is maximal. Eliminating the remaining positive-dimensional fibers in every dimension is not accomplished here. The general problem remains unresolved by this attempt.

\subsection{Sources}
\begin{itemize}
\item[S1]On the Greenfield-Wallach and Katok conjectures.
\url{https://arxiv.org/pdf/0706.3981}.
\textit{Formulation source.}
\end{itemize}
\end{document}
